===========================================================================================================================================================================================
===========================================================================================================================================================================================
===========================================================================================================================================================================================
Vektpri (kasnija cjelina)
===========================================================================================================================================================================================
===========================================================================================================================================================================================
===========================================================================================================================================================================================
Čelcoka 1
===========================================================================================================================================================================================
===========================================================================================================================================================================================
===========================================================================================================================================================================================

===========================================================================================================================================================================================
===========================================================================================================================================================================================
===========================================================================================================================================================================================
Čelcoka 2
===========================================================================================================================================================================================
===========================================================================================================================================================================================
===========================================================================================================================================================================================

===========================================================================================================================================================================================
===========================================================================================================================================================================================
===========================================================================================================================================================================================
Čelcoka 3
===========================================================================================================================================================================================
===========================================================================================================================================================================================
===========================================================================================================================================================================================

===========================================================================================================================================================================================
===========================================================================================================================================================================================
===========================================================================================================================================================================================
Čelcoka 4
===========================================================================================================================================================================================
===========================================================================================================================================================================================
===========================================================================================================================================================================================
Vektori i kolotur
Uloga kolotura ( čekrka) u dizanju tereta poznata je od davnina. Važno je dobro razumjeti principe njegovog rada koji se temelje na rastavljanju vektora u komponente.
Na prvoj slici prikazana je statična situacija kolotura u ravnoteži s prikazom sila koje djeluju na užad i na kolo- ture, te utezima koji te koloture drže u ravnoteži. Kolotur mijenja silu po njezinom smjeru, ali čuva iznos sile. Sila koja djeluje na pojedini kolotur, a koja je kom- penzirana time što je kolotur učvršćen u zid, jednaka je zbroju sila koje djeluju na užad oko tog kolotura. U elektroničkom dijelu udžbenika priložen je programčić koji opisuje ovu situaciju. Pomicanjem hvatišta tereta mijenjaju se i komponente sile u smjerovima užadi.Tu je važno shvatiti što se doga da - kolotura postaje ravno. Komponente sile tad se povećavaju preko svake granice. To objašnjava uže izme du - dva zida tako da bude potpuno ravno. Njegova težina predstavlja zašto uže nikad ne možemo razvući izme du uteg obješen o to uže i kod ravnog bi užeta djelovale na njegove krajeve beskonačno velike sile. Dakle, prije nego ga potpuno napnemo, ili bi puknulo uže, ili bi izletio vijak iz zida. Na slici iznad pokazana je uloga višestrukih kolotura u smanjenju sile pri podizanju tereta. U priloženom prog- ramčiću dana je ilustracija s četiri kolotura pri čemu se sila potrebna za podizanje tereta smanjuje četiri puta. Cijena koju treba pritom platiti je četverostruko dulje uže koje treba povući u odnosu na visinu na koju podignemo teret. Moderni kranovi koriste konstrukcije s desetak, pa i više takvih kolotura. Ovaj je princip poznat odavna. Potraži na internetu slike i objašnjenja rada dvostrukih i višes- trukih kolotura koje se još od Arhimedovih vremena rabe naročito u graditeljstvu i pomorstvu. Neke od tih slika vidimo dolje.

===========================================================================================================================================================================================
===========================================================================================================================================================================================
===========================================================================================================================================================================================
Čelcoka 5
===========================================================================================================================================================================================
===========================================================================================================================================================================================
===========================================================================================================================================================================================

===========================================================================================================================================================================================
===========================================================================================================================================================================================
===========================================================================================================================================================================================
Čelcoka 6
===========================================================================================================================================================================================
===========================================================================================================================================================================================
===========================================================================================================================================================================================

===========================================================================================================================================================================================
===========================================================================================================================================================================================
===========================================================================================================================================================================================
Čelcoka 7
===========================================================================================================================================================================================
===========================================================================================================================================================================================
===========================================================================================================================================================================================
Ako je u kvadratnoj mreži smješten mnogokut čiji su svi vrhovi u točkama --- sjecištima pravaca te mreže, tada se površina tog mnogokuta može izračunati po Pickovoj formuli1. Neka je u broj točaka unutar mnogokuta, a r broj točaka na rubu mnogokuta. Tada je površina mnogokuta jednakaP=u+12r−1kvadratnih jedinica. Kvadratna jedinica je površina najmanjeg kvadrata mreže.
Površina sedmerokuta na slici je jednaka P=5+112−1=9.5 kv. jed.
Na slici su nacrtana još tri mnogokuta smještena tako da su im svi vrhovi u točkama kvadratične mreže. Izračunajmo im površinu.
Za prvi od njih, kvadrat, imamo da je u=5, r=8 te je P=5+12⋅8−1=8.
Možemo se uvjeriti u točnost rezultata. Povučemo li dijagonale kvadrata dobit ćemo četiri sukladna pravokutna trokuta koji se mogu složiti u dva kvadrata, svaki sa stranicom duljine 2. Tako je površina jednaka 2⋅4=8 kv. jed.
Površina trokuta iznosi P=3+12⋅5−1=92, a površina šesterokuta je jednaka P=2+12⋅5−1=72=3.5 kv. jed.
Koristeći se Pickovom formulom možemo približno odrediti površinu nekog područja na zemljopisnoj karti. Na prozirni papir nacrtamo kvadratičnu mrežu te njome pokrijemo područje kojem želimo odrediti površinu. Zatim na papiru iscrtamo mnogokut s vrhovima u točkama mreže nastojeći što bolje pokriti to područje.
Slijedi određivanje brojeva u i r te računanje površine Pickovom formulom.
Dodajmo još kako je ova lijepa formula ekvivalentna čuvenoj Eulerovoj formuli za poliedre, o kojoj smo učili u II. razredu.
1 Georg Alexander Pick, austrijski matematičar, rođen u Beču 1859., umro (najvjerojatnije 1943.) u nacističkom koncentracijskom logoru u Theresienstadtu. Svoj teorem, o kojem je ovdje riječ, objavio je 1899. godine.
